% =============================================================================
% 6. RELATED WORK
% =============================================================================
\section{Related Work}
\vspace{-0.5ex}
\label{sec:related}

\textbf{Mid-training and continued pretraining.}
A growing line of work identifies a distinct stage between large-scale
pretraining and task-specific post-training, in which a model is further
trained on a smaller, carefully curated corpus to install inductive biases
that are hard to acquire from generic web text or downstream fine-tuning.
MiniCPM~\citep{minicpm} and OLMo~\citep{olmo} schedule such a stage to
upweight high-quality or domain-relevant data near the end of pretraining,
and Code-Llama~\citep{codellama} long-context adaptation can be read as a
mid-training stage targeting context-length generalization. We adopt the
same staging philosophy but apply it to a different end: installing an
\textit{agent-oriented} structural prior through function-aware FIM.

\textbf{Fill-in-the-middle objectives.}
The fill-in-the-middle objective~\citep{bavarian2022fim} is a standard
ingredient in modern code LLM pretraining, with random-span variants used
in Code-Llama~\citep{codellama}, StarCoder~\citep{starcoder},
Qwen-Coder~\citep{qwencoder}, and DeepSeek-Coder~\citep{deepseekcoder}.
Our recipe extends and specializes prior FIM along three axes. First,
masking targets are chosen at \textit{function granularity} via program
dependency graph analysis and a complexity--inferability double criterion,
so the masked region corresponds to a unit of reasoning rather than an
arbitrary span. Second, we embed an \textit{explicit chain-of-thought}
inside the FIM middle, supervising the model to produce a rationale
before the code. Third, the objective lives in a \textit{dedicated
mid-training stage} rather than dissolved into pretraining.

\textbf{Coding agent foundation models.}
Progress on coding agents has been driven by benchmarks such as
SWE-Bench~\citep{swebench}, scaffolds such as SWE-agent~\citep{sweagent}
and OpenHands~\citep{openhands}, and trajectory-centric post-training
pipelines such as SWE-Gym~\citep{swegym}, r2e-gym~\citep{r2egym}, and
SWE-Smith~\citep{swesmith}. Our work is complementary: we use r2e-gym and
SWE-Smith unmodified on top of a mid-trained checkpoint, and the novelty
lies one stage earlier---extracting a self-supervised signal from the
structure already present in source code, which both raises in-domain
agent metrics and mitigates the off-domain capability erosion that
trajectory-only post-training otherwise inflicts.

\textbf{Distillation from frontier models.}
Our default recipe uses Gemini-3-Preview to generate the chain-of-thought
rationales embedded in each FIM middle span, placing it within the
tradition of distilling capability from a strong teacher~\citep{orca,
selfinstruct, distillwhisper}. We isolate the contribution of FIM
structure from CoT distillation through a controlled ablation in
Section~\ref{sec:exp:ablation}: removing the rationale still beats the
no-mid-training baseline, and replacing the frontier teacher with
self-generated rationales recovers a sizeable fraction of the gain, so
the recipe is not distillation-bound.
